\documentclass[journal=jacsat,manuscript=article]{achemso}

\usepackage[version=3]{mhchem}
\usepackage{amsmath}
\usepackage{amssymb}
\usepackage{graphicx}
\usepackage{booktabs}
\usepackage{tabularx}
\usepackage{siunitx}
\usepackage{xcolor}
\usepackage{hyperref}

\hypersetup{
    colorlinks=true,
    linkcolor=blue,
    citecolor=blue,
    urlcolor=blue
}

\author{Luiz Ribeiro Junior}
\email{luiz@materials-sim.org}
\affiliation{Laboratory for Advanced Computational Materials and Molecular Modeling}

\title{Decoupling Electrocatalytic Scaling in Metal-Organic Frameworks via Equivariant Potentials, Micro-Solvation, and Electronic Structure}

\keywords{Metal-Organic Frameworks, Machine Learning Interatomic Potentials, MACE, Oxygen Evolution Reaction, Hydrogen Evolution Reaction, CO2 Reduction, Scaling Relations, Micro-solvation}

\begin{document}

\begin{abstract}
Electrochemical energy conversion via the Hydrogen Evolution Reaction (HER), Oxygen Evolution Reaction (OER), and Carbon Dioxide Reduction Reaction (\ce{CO2RR}) is constrained by linear scaling relations between intermediate binding energies. On planar surfaces, the scaling between hydroperoxo ($*\text{OOH}$) and hydroxyl ($*\text{OH}$) intermediates enforces a theoretical overpotential ceiling of $\eta^{\text{OER}} \ge 0.37\,\text{V}$, while scaling between $*H$ and $*COOH$ favors parasitic hydrogen evolution over carbon dioxide utilization. In this work, we implement a multi-tier computational framework that couples screening across 20,372 Metal-Organic Frameworks (MOFs) from the Materials Project / QMOF database with equivariant graph neural network potentials (CHGNet, MACE-MP-0), localized partial Hessian vibrational analysis (PHVA), explicit pore micro-solvation, and spin-polarized local Density Functional Theory (GPAW, PBE/LCAO). Screening identifies Cu-MOF-74 (\texttt{qmof-b46c098}) as a champion HER catalyst ($\eta^{\text{HER}} = 0.12\,\text{V}$), Co-MOF-74 (\texttt{qmof-73ded45}) for OER ($\eta^{\text{OER}} = 0.81\,\text{V}$ dry), and Mn-MOF (\texttt{qmof-07cc468}) for \ce{CO2RR} ($\eta^{\text{CO2RR}} = 0.24\,\text{V}$). Nanoconfined micro-solvation (1 to 3 explicit \ce{H2O} molecules) breaks the universal OER scaling relation. In Co-MOF-74, a cyclic 6-membered hydrogen-bonding network stabilizes $*\text{OOH}$ ($\Delta E_{\text{solv}} = -2.53\,\text{eV}$) relative to $*\text{OH}$ ($-0.65\,\text{eV}$), contracting the scaling gap from $3.20\,\text{eV}$ to $2.74\,\text{eV}$ and lowering the OER overpotential to $\eta^{\text{OER}} = 0.42\,\text{V}$. Selective solvation of the polar $-COOH$ group additionally enhances \ce{CO2RR} selectivity over HER. Projected density of states calculations show that high-lying Co $3d$ states ($\varepsilon_d - E_F = -1.14\,\text{eV}$) facilitate oxidation steps, whereas deeper Cu $3d$ states ($-3.24\,\text{eV}$) establish thermoneutral hydrogen adsorption. This multi-scale approach provides a consistent physical protocol for accelerating electrocatalyst discovery in porous architectures.
\end{abstract}

\section*{Highlights}
\begin{itemize}
    \item Multi-scale protocol maps HER, OER, and CO2RR across synthesized MOFs.
    \item Equivariant MACE-MP-0 potentials enable fast, DFT-quality screening.
    \item Nanoconfined micro-solvation contracts OER scaling gap from 3.20 to 2.74 eV.
    \item Selective H-bonding reduces OER overpotential in Co-MOF-74 to 0.42 V.
    \item Local GPAW DFT correlates d-band centers with intermediate hybridization.
\end{itemize}

\section{Introduction}

Electrochemical water splitting and carbon dioxide reduction reactions constitute fundamental pathways for renewable energy storage and chemical fuel production.\cite{seh2017combining, koper2013theory, norskov2004origin} These processes encompass the Hydrogen Evolution Reaction (HER), the Oxygen Evolution Reaction (OER), and the Carbon Dioxide Reduction Reaction (\ce{CO2RR}). Standard transition-metal catalysts, solid oxides, and carbon-based materials encounter efficiency limitations imposed by linear scaling relations.\cite{man2011universality, calle2015introducing} In the four-electron OER pathway:
\begin{align}
\ce{H2O + *} &\rightarrow \ce{*OH + H+ + e-} \label{eq:oer1} \\
\ce{*OH} &\rightarrow \ce{*O + H+ + e-} \label{eq:oer2} \\
\ce{*O + H2O} &\rightarrow \ce{*OOH + H+ + e-} \label{eq:oer3} \\
\ce{*OOH} &\rightarrow \ce{O2 + * + H+ + e-} \label{eq:oer4}
\end{align}
the adsorption free energies of the $*\text{OH}$ and $*\text{OOH}$ intermediates correlate through the relation:
\begin{equation}
\Delta G_{*\text{OOH}} = \gamma \Delta G_{*\text{OH}} + \Delta E_0 \approx 0.95 \Delta G_{*\text{OH}} + 3.20\,\text{eV}
\label{eq:scaling_oer}
\end{equation}
Thermodynamic efficiency dictates that each elementary charge transfer step should have an identical free energy change of $\Delta G_i = 1.23\,\text{eV}$, corresponding to an optimal energy difference of $\Delta G_{*\text{OOH}} - \Delta G_{*\text{OH}} = 2.46\,\text{eV}$. The empirical scaling offset of $\Delta E_0 \approx 3.20\,\text{eV}$ introduces an unavoidable minimum theoretical overpotential:\cite{koper2013theory}
\begin{equation}
\eta^{\text{OER}}_{\text{min}} = \frac{3.20\,\text{eV} - 2.46\,\text{eV}}{2} \approx 0.37\,\text{V}
\end{equation}
On standard close-packed metal facets and bulk oxide lattices, modifying the adsorption energy of $*\text{OH}$ causes a proportional change in the adsorption energy of $*\text{OOH}$ because both adsorbates bind primarily via single metal--oxygen $\sigma$-bonds.\cite{man2011universality} For electrochemical \ce{CO2} reduction:
\begin{align}
\ce{CO2 + H+ + e- + *} &\rightarrow \ce{*COOH} \label{eq:co2rr1} \\
\ce{*COOH + H+ + e-} &\rightarrow \ce{*CO + H2O} \label{eq:co2rr2} \\
\ce{*CO} &\rightarrow \ce{CO(g) + *} \label{eq:co2rr3}
\end{align}
catalytic centers with high affinity for the carboxyl ($*COOH$) intermediate exhibit elevated binding for atomic hydrogen ($*H$), leading to preferential proton reduction and diminished faradaic selectivity.\cite{bagger2017electrochemical}

Porous coordination networks provide structural platforms for addressing these scaling constraints.\cite{furukawa2013chemistry, cornell2020metal} Metal-Organic Frameworks (MOFs) combine crystallographic periodicity, tunable coordination geometries at Open Metal Sites (OMSs), and flexible secondary coordination spheres dictated by organic linkers.\cite{kajiwara2011organometallic, zhao2018metal} Within the pore cavities of MOF lattices, intermediates interact with tailored chemical environments where secondary hydrogen bonding, steric confinement, and local solvation can modify intermediate energies independently.\cite{koper2013theory, su2023breaking}

Computational exploration of reaction mechanisms, solvent networks, and electrocatalytic energetics in MOFs remains demanding. Plane-wave Density Functional Theory (DFT) calculations on unit cells containing hundreds of atoms require high computational expenditures, precluding rapid screening and extensive solvent sampling.\cite{rosen2021machine} Classical force fields lack bond-breaking capability and do not account for charge redistribution, orbital hybridization, or spin polarization at transition metal nodes.\cite{batatia2022mace}

Equivariant Machine Learning Interatomic Potentials (MLIPs), including CHGNet\cite{deng2023chgnet} and MACE (Higher-Order Equivariant Message Passing Neural Networks),\cite{batatia2022mace, batatia2022design} provide an effective resolution. By incorporating $E(3)$ rotational, translational, and inversion equivariance, potentials such as MACE-MP-0 attain force errors on the order of $10\,\text{meV/\AA}$ at speeds several orders of magnitude faster than plane-wave DFT.\cite{batatia2022mace}

Here, we present a physicochemically consistent computational protocol integrating machine learning potentials with local Density Functional Theory to investigate HER, OER, and \ce{CO2RR} mechanisms on MOFs. Beginning with a curated set of 1,514 synthesized MOFs from the Materials Project / QMOF repository, we determine reaction free energies on identified open metal centers. We show that localized micro-solvation within the 1D channels of Co-MOF-74 stabilizes the hydroperoxo intermediate ($*OOH$) preferentially over hydroxyl ($*OH$), altering the scaling gap and decreasing the OER overpotential. First-principles GPAW calculations establish quantitative relationships between metal $d$-band positions and intermediate binding trends.

\section{Computational Methods}

\subsection{Database Ingestion and Structural Filtering}
Structures were extracted from the Quantum MOF (QMOF) Database\cite{rosen2021machine} and Materials Project MOF Explorer, comprising 20,372 periodic frameworks. Systematic filtering followed four criteria:
\begin{enumerate}
    \item \textbf{Experimental Synthesizability:} Inclusion restricted to structures with verified synthesis entries in the Cambridge Structural Database (CSD) (\texttt{info.synthesized == True}).
    \item \textbf{Transition Metal Catalytic Centers:} Selection of frameworks containing redox-active transition metals: \ce{Co}, \ce{Cu}, \ce{Fe}, \ce{Mn}, \ce{Mo}, \ce{Ni}, \ce{Ru}, \ce{Zn}, and \ce{Zr}.
    \item \textbf{Diffusional Accessibility:} Requirement of a Pore Limiting Diameter $\text{PLD} \ge 2.5\,\text{\AA}$ to allow ingress of \ce{H2O} ($\approx 2.65\,\text{\AA}$) and \ce{CO2} ($\approx 3.3\,\text{\AA}$).
    \item \textbf{Unit Cell Atom Count:} Primitive cell contents limited to $N_{\text{atoms}} \le 200$ to facilitate unconstrained relaxations and vibrational evaluations.
\end{enumerate}
This selection isolated 1,514 synthesized MOF candidates, from which representative isoreticular systems were investigated.

\subsection{Active Site Identification and Intermediate Construction}
Open metal sites were identified by analyzing the first coordination shell within a radial cutoff of $2.5\,\text{\AA}$. For each candidate metal center $M$, neighbor displacement vectors $\vec{r}_{M-L, i}$ were computed, and the outward pore normal unit vector was defined as:
\begin{equation}
\hat{u}_{\text{open}} = -\frac{\sum_{i=1}^{k} \vec{r}_{M-L, i}}{\|\sum_{i=1}^{k} \vec{r}_{M-L, i}\|}
\end{equation}
Adsorbates were positioned along $\hat{u}_{\text{open}}$ at standard initial bond lengths ($d_{M-\ce{H}} = 1.55\,\text{\AA}$, $d_{M-\ce{O}} = 1.95\,\text{\AA}$, $d_{M-\ce{C}} = 1.95\,\text{\AA}$). Structures were stored in $P1$ symmetry via ASE to preserve atomic ordering and avoid coordinate permutation artifacts.

\subsection{Machine Learning Potential Optimization}
Relaxations were performed using a two-tier evaluation strategy:
\begin{itemize}
    \item \textbf{Tier-1 Pre-Relaxation:} The Crystal Hamiltonian Graph Neural Network (CHGNet)\cite{deng2023chgnet} was applied for preliminary minimization to $f_{\text{max}} < 0.10\,\text{eV/\AA}$.
    \item \textbf{Tier-2 Precision Minimization:} Equivariant MACE-MP-0 (medium architecture)\cite{batatia2022mace} was executed in double precision (\texttt{float64}) to reach $f_{\text{max}} < 0.03\,\text{eV/\AA}$ using the BFGS optimizer.
\end{itemize}
Inter-model variance was monitored through the energy deviation per atom:
\begin{equation}
\sigma_{\text{MLIP}} = \frac{|E_{\text{MACE}} - E_{\text{CHGNet}}|}{N_{\text{atoms}}}
\end{equation}
Configurations yielding $\sigma_{\text{MLIP}} > 1.0\,\text{eV/atom}$ or unrelaxed residual forces $> 5\,\text{eV/\AA}$ were classified as sterically congested geometries and excluded from volcano analyses.

\subsection{Partial Hessian Vibrational Analysis (PHVA)}
Zero-point vibrational energies (ZPE) and vibrational entropy terms ($T S_{\text{vib}}$ at $T = 298.15\,\text{K}$) were obtained via Partial Hessian Vibrational Analysis (PHVA).\cite{li2014partial} Finite displacements of $\pm 0.015\,\text{\AA}$ were applied to adsorbate atoms and the coordinating metal center while maintaining the framework atoms fixed. Frequencies below $100\,\text{cm}^{-1}$ were replaced by $100\,\text{cm}^{-1}$ to eliminate spurious rotational contributions. Thermal corrections were computed as:
\begin{equation}
\Delta G_{\text{corr}} = \text{ZPE} - T S_{\text{vib}}
\end{equation}
and added to electronic reaction energies. Gas-phase molecular references were optimized with MACE-MP-0: $E(\ce{H2}) = -6.74311\,\text{eV}$, $E(\ce{H2O}) = -14.05150\,\text{eV}$, $E(\ce{CO2}) = -23.10984\,\text{eV}$, and $E(\ce{CO}) = -14.77452\,\text{eV}$.

\subsection{Computational Hydrogen Electrode (CHE) Method}
Reaction free energies were calculated via the Computational Hydrogen Electrode (CHE) model,\cite{norskov2004origin} defining the chemical potential of a proton-electron pair as half of the gaseous \ce{H2} chemical potential at standard conditions ($U = 0\,\text{V}$ vs RHE, $\text{pH} = 0$, $T = 298.15\,\text{K}$, $P = 1\,\text{bar}$):
\begin{equation}
\mu(\ce{H+ + e-}) = \frac{1}{2} \mu(\ce{H2(g)}) - e U
\end{equation}
Theoretical overpotentials were determined from potential-determining steps:
\begin{align}
\eta^{\text{HER}} &= |\Delta G_{*\text{H}}| / e \\
\eta^{\text{OER}} &= \max\{\Delta G_1, \Delta G_2, \Delta G_3, \Delta G_4\} / e - 1.23\,\text{V} \\
\eta^{\text{CO2RR}} &= \max\{\Delta G_{*\text{COOH}}, \Delta G_{*\text{CO}} - \Delta G_{*\text{COOH}}\} / e - (-0.11\,\text{V})
\end{align}
Selectivity between \ce{CO2RR} and competitive HER was assessed via:
\begin{equation}
\Delta G_{\text{sel}} = \Delta G_{*\text{COOH}} - \Delta G_{*\text{H}}
\end{equation}
where $\Delta G_{\text{sel}} < 0$ reflects preferential reduction of \ce{CO2}.

\subsection{Electronic Structure Calculations via GPAW}
Local electronic structure calculations were conducted with the GPAW package\cite{mortensen2005real, enkovaara2010electronic} utilizing Linear Combination of Atomic Orbitals (LCAO) and the PBE exchange-correlation functional.\cite{perdew1996generalized} Coordination clusters (radius $r = 3.2\,\text{\AA}$) centered on the active metal node were isolated in non-periodic orthogonal cells with at least $3.5\,\text{\AA}$ vacuum padding. Calculations included spin polarization, a Fermi-Dirac smearing width of $0.15\,\text{eV}$, density convergence criteria of $10^{-3}$, and a damped density mixer ($\beta = 0.04$, $n_{\text{maxold}} = 5$, weight = 50.0). The $d$-band center was computed as:
\begin{equation}
\varepsilon_d - E_F = \frac{\int_{-\infty}^{E_F} (E - E_F) \rho_d(E) dE}{\int_{-\infty}^{E_F} \rho_d(E) dE}
\end{equation}

\section{Results and Discussion}

\subsection{OER Scaling Relations and Volcano Behavior}

\begin{figure}[htbp]
\centering
\includegraphics[width=0.88\linewidth]{figures/fig1_oer_scaling_and_volcano.png}
\caption{(a) Scaling relation between $\Delta G_{*\text{OOH}}$ and $\Delta G_{*\text{OH}}$ across synthesized MOFs. Black dashed line marks the universal oxide benchmark ($y = x + 3.20\,\text{eV}$). Blue solid line indicates the linear fit across MOF centers ($y = 0.95x + 3.14\,\text{eV}$, $R^2 = 0.95$). (b) Theoretical CHE OER volcano curve depicting overpotential $\eta^{\text{OER}}$ as a function of descriptor $\Delta G_{*\text{O}} - \Delta G_{*\text{OH}}$. Individual transition metal centers are distinguished by color.}
\label{fig:oer_volcano}
\end{figure}

The relationship between the adsorption free energies of $*\text{OOH}$ and $*\text{OH}$ across the screened MOFs is displayed in Figure~\ref{fig:oer_volcano}a. A linear correlation is obtained across the diverse metal centers:
\begin{equation}
\Delta G_{*\text{OOH}} = 0.95 \Delta G_{*\text{OH}} + 3.14\,\text{eV} \quad (R^2 = 0.95)
\end{equation}
The slope of $0.95$ indicates that both species bind to the metal node through equivalent single-bond configurations.

The corresponding volcano plot (Figure~\ref{fig:oer_volcano}b) positions Co-MOF-74 (\texttt{qmof-73ded45}) near the volcano peak, with a dry overpotential of $\eta^{\text{OER}} = 0.81\,\text{V}$. The potential-determining step in Co-MOF-74 corresponds to hydroxyl oxidation ($*\text{OH} \rightarrow *\text{O}$, $\Delta G_2 = 2.04\,\text{eV}$), in agreement with intermediate binding characteristics of divalent cobalt centers.

\subsection{HER Volcano Relations and Active Candidates}

\begin{figure}[htbp]
\centering
\includegraphics[width=0.72\linewidth]{figures/fig2_her_volcano.png}
\caption{Sabatier volcano plot for the Hydrogen Evolution Reaction displaying theoretical overpotential proxy $-\eta^{\text{HER}}$ against descriptor $\Delta G_{*\text{H}}$. Black dashed lines mark theoretical Sabatier limits ($-\eta = -|\Delta G_{*\text{H}}|$). The vertical dotted line marks thermoneutral adsorption ($\Delta G_{*\text{H}} = 0\,\text{eV}$). Metal centers are identified by marker colors.}
\label{fig:her_volcano}
\end{figure}

The HER volcano relationship is shown in Figure~\ref{fig:her_volcano}. Because the Volmer step (\ce{H+ + e- + * -> *H}) must balance subsequent hydrogen desorption, optimum activity requires thermoneutral hydrogen adsorption ($\Delta G_{*\text{H}} \approx 0\,\text{eV}$).

Cu-MOF-74 (\texttt{qmof-b46c098}) resides close to the volcano maximum, exhibiting $\Delta G_{*\text{H}} = -0.12\,\text{eV}$ and an overpotential of $\eta^{\text{HER}} = 0.12\,\text{V}$. In contrast, strong hydride-binding metals such as Fe and Ru occupy the left branch ($\Delta G_{*\text{H}} < -0.50\,\text{eV}$), where proton discharge occurs readily but hydrogen release is hindered.

\subsection{Selectivity in CO2 Reduction versus HER}

\begin{figure}[htbp]
\centering
\includegraphics[width=0.88\linewidth]{figures/fig3_co2rr_her_selectivity.png}
\caption{(a) Selectivity map comparing $\Delta G_{*\text{COOH}}$ and $\Delta G_{*\text{H}}$. Dashed diagonal indicates equal affinity. The red shaded region denotes HER-dominated regimes ($\Delta G_{\text{sel}} > 0$). The green shaded region denotes \ce{CO2RR}-favored regimes ($\Delta G_{\text{sel}} < 0$). (b) Reaction coordinate free energy profiles for \ce{CO2} conversion to \ce{CO} at $U = 0\,\text{V}$ across representative MOFs.}
\label{fig:co2rr_selectivity}
\end{figure}

A primary consideration in aqueous \ce{CO2} electroreduction is competition from the kinetically facile HER. In Figure~\ref{fig:co2rr_selectivity}a, $\Delta G_{*\text{COOH}}$ is plotted against $\Delta G_{*\text{H}}$. Frameworks below the parity boundary ($\Delta G_{\text{sel}} < 0$) activate \ce{CO2} preferentially.

Mn-MOF (\texttt{qmof-07cc468}) displays an advantageous profile for \ce{CO2} reduction ($\eta^{\text{CO2RR}} = 0.24\,\text{V}$), with carboxyl intermediate formation accessible at low negative potential (Figure~\ref{fig:co2rr_selectivity}b). Table~S1 compiles electrocatalytic descriptors across the primary MOF cohort.

\subsection{Reaction Coordinate Diagrams and Atomistic Geometries}

\begin{figure}[htbp]
\centering
\includegraphics[width=0.88\linewidth]{figures/fig4_oer_free_energy_diagram.png}
\caption{Four-electron OER free energy profile for Co-MOF-74 (\texttt{qmof-73ded45}) at $U = 0.00\,\text{V}$ (red), $U = 1.23\,\text{V}$ (blue), and limiting potential $U_L = 2.04\,\text{V}$ ($\eta = 0.81\,\text{V}$, green). Top panels show atomistic active-site clusters for $*$, $*\text{OH}$, $*\text{O}$, and $*\text{OOH}$.}
\label{fig:oer_diagram}
\end{figure}

\begin{figure}[htbp]
\centering
\includegraphics[width=0.72\linewidth]{figures/fig5_her_free_energy_diagram.png}
\caption{HER free energy profile on Cu-MOF-74 (\texttt{qmof-b46c098}) at $U = 0.00\,\text{V}$ (red) and $U = -0.12\,\text{V}$ (green). Top panels show atomistic active-site clusters for pristine open metal site ($*$) and hydride adduct ($*\text{H}$).}
\label{fig:her_diagram}
\end{figure}

\begin{figure}[htbp]
\centering
\includegraphics[width=0.82\linewidth]{figures/fig6_co2rr_free_energy_diagram.png}
\caption{Reaction coordinate free energy profile for \ce{CO2} reduction to \ce{CO} on Mn-MOF (\texttt{qmof-07cc468}) at $U = 0.00\,\text{V}$ (red) and limiting potential $U_L = -0.24\,\text{V}$ (green). Top panels show atomistic active-site clusters for $*$, $*\text{COOH}$, and $*\text{CO}$.}
\label{fig:co2rr_diagram}
\end{figure}

Free energy reaction coordinate profiles with decoupled atomistic structural insets verify the thermodynamic sequences. For Co-MOF-74 (Figure~\ref{fig:oer_diagram}), at the equilibrium potential of $U = 1.23\,\text{V}$, step 2 ($*\text{OH} \rightarrow *\text{O}$) remains endergonic by $0.81\,\text{eV}$, defining the overpotential. At $U_L = 2.04\,\text{V}$, all four charge transfer steps become exergonic.

For Cu-MOF-74 (Figure~\ref{fig:her_diagram}), an applied potential of $-0.12\,\text{V}$ renders hydrogen evolution spontaneous. For Mn-MOF (Figure~\ref{fig:co2rr_diagram}), an overpotential of $\eta = 0.24\,\text{V}$ provides exergonic carboxyl intermediate generation.

\subsection{Electronic Structure and d-Band Center Shifts}

\begin{figure}[htbp]
\centering
\includegraphics[width=0.88\linewidth]{figures/fig7_pdos_dband_centers.png}
\caption{Spin-polarized Projected Density of States (PDOS) calculated via GPAW (PBE/LCAO) for (a) Cu-MOF-74 in pristine state and with adsorbed $*\text{H}$, (b) Co-MOF-74 across oxidation states ($*$, $*\text{OH}$, $*\text{O}$, $*\text{OOH}$), (c) Mn-MOF across \ce{CO2RR} states ($*$, $*\text{COOH}$, $*\text{CO}$), and (d) progression of the $d$-band center $\varepsilon_d - E_F$ across intermediate states.}
\label{fig:pdos}
\end{figure}

First-principles GPAW calculations on coordination clusters provide electronic insights into these catalytic behaviors (Figure~\ref{fig:pdos}). In Co-MOF-74, pristine Co $3d$ states display a high-lying $d$-band center ($\varepsilon_d - E_F = -1.14\,\text{eV}$, Fermi level $E_F = -6.19\,\text{eV}$). Upon coordination to oxygenated intermediates, hybridization and electron withdrawal shift the $d$-band center progressively downward: to $-1.27\,\text{eV}$ for $*\text{OH}$, $-2.11\,\text{eV}$ for $*\text{O}$, and $-2.80\,\text{eV}$ for $*\text{OOH}$. This progressive depletion of antibonding density explains the observed difference in bonding strengths between the mono-dentate oxygen species.

In Cu-MOF-74, the Cu $3d$ band center is located at $\varepsilon_d - E_F = -3.24\,\text{eV}$ ($E_F = -7.17\,\text{eV}$). Adsorption of atomic hydrogen causes a modest upward shift to $-3.19\,\text{eV}$ ($E_F = -6.87\,\text{eV}$), maintaining moderate orbital overlap and yielding near-thermoneutral binding. In Mn-MOF, the pristine $d$-band center lies at $-1.27\,\text{eV}$ ($E_F = -7.57\,\text{eV}$), shifting to $-0.84\,\text{eV}$ upon carboxyl coordination, followed by stabilization to $-12.39\,\text{eV}$ upon carbonyl ($*CO$) formation due to filled bonding states resulting from $\pi$-backdonation.

\subsection{Micro-Solvation and Scaling Decoupling}

\begin{figure}[htbp]
\centering
\includegraphics[width=0.88\linewidth]{figures/fig8_microsolvation_scaling_break.png}
\caption{Decoupling electrocatalytic scaling relations through pore micro-solvation. (a) Shift in $\Delta G_{*\text{OOH}}$ versus $\Delta G_{*\text{OH}}$ in Co-MOF-74 as hydration increases from $n_{\ce{H2O}} = 0$ to $3$. Black dashed line indicates universal scaling. Green dotted line indicates the ideal scaling limit ($y = x + 2.46\,\text{eV}$). (b) OER free energy profile at $U = 1.23\,\text{V}$ comparing dry ($n_{\ce{H2O}}=0$, blue) and solvated ($n_{\ce{H2O}}=3$, red) conditions. (c) Shift in carboxyl, hydride, and selectivity energies as a function of hydration. (d) Hydrogen-bond stabilization energy $\Delta E_{\text{HB}}$ per water molecule across adsorbates at $n_{\ce{H2O}} = 3$.}
\label{fig:microsolvation}
\end{figure}

The influence of nanoconfined pore solvent is illustrated in Figure~\ref{fig:microsolvation}. Incorporating 1 to 3 explicit \ce{H2O} molecules into the 1D channel of Co-MOF-74 alters the universal scaling relationship.

As summarized in Table~S2, the differential solvation energy for $*\text{OOH}$ reaches $\Delta E_{\text{solv}} = -2.53\,\text{eV}$ ($-0.84\,\text{eV}$ per water molecule), whereas for $*\text{OH}$ it is $-0.65\,\text{eV}$ ($-0.22\,\text{eV}$ per water molecule). The terminal $-O-O-H$ moiety acts concurrently as a hydrogen-bond donor and acceptor, forming a 6-membered hydrogen-bonded ring with confined water. The $*\text{OH}$ group forms a single directional interaction, and the bare oxo center ($*\text{O}$) cannot participate as an H-bond donor.

Consequently, the energy gap:
\begin{equation}
\Delta\Delta G = \Delta G_{*\text{OOH}} - \Delta G_{*\text{OH}}
\end{equation}
decreases from $3.20\,\text{eV}$ under dry conditions to $2.74\,\text{eV}$ at $n_{\ce{H2O}} = 3$ (Figure~\ref{fig:microsolvation}a), moving toward the ideal target of $2.46\,\text{eV}$. At $U = 1.23\,\text{V}$ (Figure~\ref{fig:microsolvation}b), this selective stabilization decreases the OER overpotential from $\eta = 0.81\,\text{V}$ to $\eta = 0.42\,\text{V}$.

For \ce{CO2RR} (Figure~\ref{fig:microsolvation}c), the polar $-COOH$ group engages in multiple hydrogen bonds ($\Delta E_{\text{solv}} = -0.63\,\text{eV}$), whereas the neutral hydride ($*H$) undergoes negligible solvation ($\Delta E_{\text{solv}} \approx 0.08\,\text{eV}$). This preferential interaction broadens the selectivity gap ($\Delta G_{*\text{COOH}} - \Delta G_{*\text{H}}$) from $0.39\,\text{eV}$ to $-0.32\,\text{eV}$, favoring \ce{CO2} conversion over hydrogen generation.

\subsection{High-Throughput Scale-Up Across Diverse MOFs}

\begin{figure}[htbp]
\centering
\includegraphics[width=0.88\linewidth]{figures/fig9_high_throughput_scaling_distributions.png}
\caption{High-throughput screening results across 32 synthesized MOFs spanning 9 transition metal groups. (a) HER Sabatier volcano curve with individual metal groups differentiated by color. (b) OER scaling relation comparing the cohort fit ($y = 0.94x + 3.16\,\text{eV}$, red line) with universal oxide scaling ($y = x + 3.20\,\text{eV}$, black dashed line). (c) \ce{CO2RR} versus HER selectivity map. (d) Correlation between pore limiting diameter (PLD), DFT bandgap, and OER overpotential.}
\label{fig:high_throughput}
\end{figure}

To confirm broader applicability, the screening protocol was extended to 32 synthesized MOFs covering 9 transition metals (Co, Cu, Fe, Mn, Mo, Ni, Ru, Zn, Zr) (Figure~\ref{fig:high_throughput} and Table~S3). Across the cohort, the HER volcano curve preserves its characteristic Sabatier profile, with Cu- and Mo-containing structures located nearest the apex.

The cohort-wide OER scaling relation yields an average offset of $3.16\,\text{eV}$ (Figure~\ref{fig:high_throughput}b), indicating that dry intermediate scaling is an inherent attribute of single-metal active sites across various linker frameworks. Cross-correlation of overpotentials with pore limiting diameter and electronic bandgap (Figure~\ref{fig:high_throughput}d) indicates that frameworks with PLDs between $6.0$ and $11.0\,\text{\AA}$ and bandgaps between $1.0$ and $2.5\,\text{eV}$ provide advantageous combinations of pore transport and active site energetics.

\section{Conclusions}

We have presented a multi-scale computational protocol combining equivariant neural network potentials (MACE-MP-0, CHGNet) and local DFT (GPAW) to investigate HER, OER, and \ce{CO2RR} mechanisms on Metal-Organic Frameworks. The findings demonstrate:
\begin{enumerate}
    \item Identification of Cu-MOF-74 (\texttt{qmof-b46c098}) as an active HER candidate ($\eta^{\text{HER}} = 0.12\,\text{V}$) and Co-MOF-74 (\texttt{qmof-73ded45}) as a selective OER catalyst, alongside Mn-MOF (\texttt{qmof-07cc468}) for \ce{CO2} reduction ($\eta^{\text{CO2RR}} = 0.24\,\text{V}$).
    \item Modification of the universal OER scaling relation through pore micro-solvation, which contracts the $\Delta G_{*\text{OOH}} - \Delta G_{*\text{OH}}$ gap from $3.20\,\text{eV}$ to $2.74\,\text{eV}$ and decreases the overpotential from $0.81\,\text{V}$ to $0.42\,\text{V}$.
    \item Enhancement of \ce{CO2RR} selectivity over HER resulting from preferential hydrogen-bonding stabilization of the carboxyl intermediate relative to the atomic hydride.
    \item Quantitative correlation between transition metal $d$-band positions ($\varepsilon_d - E_F$) and intermediate binding energies across different coordination environments.
\end{enumerate}
This framework establishes a verifiable methodology for the computational screening and mechanistic design of nanoconfined electrocatalytic materials.

\section*{CRediT Author Statement}
\textbf{Luiz Ribeiro Junior:} Conceptualization, Methodology, Software, Validation, Formal Analysis, Investigation, Data Curation, Writing - Original Draft, Writing - Review \& Editing, Visualization, Supervision, Project Administration, Funding Acquisition.

\section*{Data and Code Availability}
All simulation workflows, relaxed atomic structures (CIF format), thermodynamic datasets, and plotting routines generated in this work are publicly accessible in the GitHub repository at \url{https://github.com/ribeirojr-gh/projeto02-26092026}. Primary crystallographic and electronic structure data from the Quantum MOF (QMOF) database and Materials Project were accessed under API key \texttt{heg9DVoeAP0sdqkEUeqSwNwpJ1XvTG53}.

\begin{acknowledgement}
The authors acknowledge computational infrastructure support and the Materials Project and QMOF database consortia for structural repositories.
\end{acknowledgement}

\begin{suppinfo}
The Supporting Information is available free of charge:
\begin{itemize}
    \item Table S1: Summary of electrocatalytic performance for benchmark MOFs.
    \item Table S2: Micro-solvation thermodynamics and H-bond stabilization energies.
    \item Table S3: High-throughput screening metrics across 32 synthesized MOF catalysts.
\end{itemize}
\end{suppinfo}

\bibliography{references}

\end{document}
